\documentclass[12pt]{article}

%% ---- Packages -------------------------------------------------------
\usepackage[margin=1in]{geometry}
\usepackage{setspace}
\usepackage{amsmath,amssymb,amsthm}
\usepackage{graphicx}
\usepackage{booktabs}
\usepackage{longtable}
\usepackage{array}
\usepackage{natbib}
\usepackage{hyperref}
\usepackage{xcolor}
\usepackage{float}
\usepackage{caption}
\usepackage{subcaption}
%\usepackage{microtype}  % disabled: font expansion not available on this install
\usepackage{enumitem}
\usepackage{threeparttable}
\usepackage{dcolumn}
\usepackage{rotating}
\usepackage{lscape}
\usepackage{multirow}

%% ---- Typography -------------------------------------------------------
\doublespacing
\setlength{\parindent}{0.5in}
\setlength{\parskip}{0pt}
\captionsetup{font=small,labelfont=bf,justification=raggedright,singlelinecheck=false}

%% ---- Math environments ------------------------------------------------
\newtheorem{proposition}{Proposition}
\newtheorem{corollary}{Corollary}
\newtheorem{lemma}{Lemma}
\newtheorem{assumption}{Assumption}
\theoremstyle{definition}
\newtheorem{definition}{Definition}
\theoremstyle{remark}
\newtheorem{remark}{Remark}

%% ---- Hyperref ----------------------------------------------------------
\hypersetup{colorlinks=true, linkcolor=blue!70!black,
            citecolor=blue!70!black, urlcolor=blue!70!black}

%% ====================================================================
\begin{document}

%% ---- Title page -------------------------------------------------------
\begin{titlepage}
\begin{center}
\vspace*{0.8cm}
{\LARGE \textbf{Rivals at the Starting Line:\\[6pt]
Peer Effects, Familiarity, and Rivalry Amplification\\[6pt]
in Competitive Track Running}}

\vspace{1.2cm}
{\large [Author(s) Redacted for Blind Review]}\\[0.4cm]
{\normalsize \today}

\vspace{1.4cm}
\begin{minipage}{0.92\textwidth}
\begin{abstract}
\noindent\singlespacing
We study how repeated competition with the same opponents reshapes the
structure and magnitude of peer effects among elite track runners.
Using 587{,}213 race observations covering 49{,}168 athletes across 5{,}180
international competitions (World Athletics Organization, 2005--2023),
we estimate all eight structural peer-effect models of \citet{HoxbyWeingarth2007}
and test how each mechanism is amplified or attenuated by the number of prior
co-races between the focal athlete and the field.
We also condition on whether repeated opponents are drawn from
high-quality or low-quality fields, separating the quality and history channels.
Four main findings emerge.
First, peer effects follow a U-shaped pattern in familiarity: large for
first-time encounters ($\hat\beta = 0.104^{***}$), attenuated after one prior
meeting ($\hat\beta = 0.032^{***}$), and sharply amplified among athletes who
have raced the same rivals seven or more times
($\hat\beta = 0.406^{***}$; a factor-of-twelve amplification).
Second, the Shining Light mechanism---whereby athletes respond to the fastest
peer in the field---dominates, while the Bad Apple effect (the slowest peer) is
absent.
Third, familiarity amplification is strongest when repeated rivals come from
high-quality fields: repeated competition against able opponents
increases the peer-effect coefficient, while repeated competition against weak
opponents attenuates it, consistent with a rivalry model in which complacency
sets in against familiar inferiors.
Fourth, familiarity amplification is most pronounced in middle-distance events
(500--2{,}000 m), where tactical pacing and rival recognition carry the highest
strategic value.
On the theory side, we extend the conformist peer-effects framework
of \citet{Boucher2024} and \citet{Bernheim1994}
by incorporating a rivalry utility motive inspired by \citet{BenabouTirole2006}
and embedding it in the tournament structure of \citet{LazearRosen1981}.
We show that all eight HOX mechanisms emerge as special cases of a general
best-reply function indexed by the choice of peer statistic and rival weighting,
derive familiarity-amplified equilibria with closed-form comparative statics,
and prove that task complexity enters through the effort cost function and
information precision, generating an inverse-U-shaped familiarity amplification
in race distance.
Importantly, because running is an individual sport with no team production,
our setting fully isolates the rivalry peer effect from the altruistic
motive---the willingness to sacrifice own performance to help teammates---that
confounds peer-effect estimates in team sports.

\vspace{0.4cm}
\noindent\textbf{JEL:} J24, Z22, D83, C23 \quad
\textbf{Keywords:} peer effects, rivalry, track and field, familiarity,
tournament, task complexity, altruism
\end{abstract}
\end{minipage}
\end{center}
\end{titlepage}

%% ====================================================================
\clearpage
\section{Introduction}
\label{sec:intro}

Athletic competition offers some of the most vivid evidence that rivals matter.
Magic Johnson reportedly began each morning by reading Larry Bird's box scores;
Roger Federer's serve reportedly stiffened whenever Rafael Nadal was on the other
side of the net; Paula Radcliffe described Gete Wami's presence as the stimulus
that took her to world-record pace in the 2003 London Marathon.
These anecdotes share a common thread: familiarity with a specific rival
intensifies motivation in a way that mere field quality does not.
The economics literature has identified many channels through which the peer
group affects individual performance---conformist preferences, tournament
incentives, information spillovers, social comparison---yet it has treated
every opponent as interchangeable.
We ask whether it matters \textit{who} you have competed against before,
not just how good they are.

This paper provides the first large-scale econometric test of whether peer effects
in a competitive environment are moderated by the history of prior interactions
between competitors.
We call this moderating variable \textit{familiarity}: the number of times
a focal athlete has previously raced the same opponent in the same event type.
Our setting is elite track and field athletics, where individual performance
is objectively measured in seconds (converted to m/s), fields are composed of
10--15 runners, opponent identities are fully observed, and prior performance
histories are available for every competitor.
These features allow us to construct all eight structural peer-effect variables
proposed by \citet{HoxbyWeingarth2007} and to condition their effects on two
dimensions---familiarity and opponent quality level---in a unified econometric
framework with athlete fixed effects and competition$\times$event fixed effects.

Our dataset draws on the global competition records of the World Athletics
Organization (WAO), covering events from 2005 to 2023 in over 1{,}500 cities
worldwide \citep{AragonRezazadeh2025}.
The panel structure---the same athletes appear in many competitions over many
years---permits construction of our key familiarity variable and credible
identification via within-athlete variation in peer composition.
The peer quality variables are constructed exclusively from opponents' \textit{prior}
performance records, predetermining them relative to the current race and
sidestepping the reflection problem \citep{Manski1993,Angrist2014}.

\paragraph{Main findings.}
We report four main results:

\begin{enumerate}[leftmargin=0.5in, label=\textit{Result \arabic*.}]

\item \textbf{U-shaped peer effects in familiarity (information shock, calibration, rivalry).}
Peer effects are largest when athletes encounter an entirely new field
($\hat\beta = 0.104$), shrink to their minimum after one prior co-race
($\hat\beta = 0.032$), and surge to more than four times the first-meeting
level when athletes have competed against the same rivals seven or more
times ($\hat\beta = 0.406$).
The drop from Never to $1\times$ reflects Bayesian updating: a first
encounter conveys high information content about the field's quality, while
the second meeting has lower marginal information value.
The subsequent surge is a genuine rivalry effect: repeated opponents are
psychologically ``owned,'' motivating effort beyond what pure information
would predict.

\item \textbf{Shining Light dominates; Bad Apple is absent.}
Among the eight HOX mechanisms, the effect of the fastest opponent's prior
speed ($\hat\beta = 0.049^{***}$) and the athlete's normalized rank in the
field ($\hat\beta = 0.076^{***}$) dominate.
The slowest opponent ($\hat\beta = -0.002$) and the Single Crossing interaction
($\hat\beta = 0.001$) are both economically negligible and statistically zero.
Boutique and Focus/Rainbow are positive but small, suggesting that pace-group
formation and field heterogeneity play secondary roles.

\item \textbf{Rivalry amplification is stronger in lower-quality fields.}
Counter-intuitively, the familiarity amplification of peer effects---the rise from
Never to $7+\times$---is larger when repeated rivals come from below-median quality fields.
In low-quality familiar fields, the Shining Light coefficient rises from 0.199 to
0.795 ($\Delta\beta = +0.596$), while in high-quality familiar fields it rises from
0.062 to 0.215 ($\Delta\beta = +0.153$).
We interpret this as a \textit{ceiling effect}: athletes in top-tier fields already
run near their personal best, leaving little room for rivalry motivation to generate
additional performance gains.
Conversely, Invidious Comparison \textit{attenuates} with familiarity in high-quality
fields (rank acceptance among familiar superiors), and the Single Crossing coefficient
reverses sign by quality stratum---positive in strong fields (peer leaders) and
negative in weak fields (underdog effect)---resolving a pooled null that
masks offsetting mechanisms.

\item \textbf{Middle-distance events exhibit the strongest familiarity amplification.}
Sprint events ($\le 400$ m) show small peer effects that barely increase with
familiarity.
Middle-distance events (500--2{,}000 m) show large peer effects that rise sharply,
reaching $\hat\beta = 0.705$ in the rivalry bin.
Long-distance events ($> 2{,}000$ m) are intermediate but suffer from thin
samples at high familiarity levels.
This pattern is consistent with our theoretical prediction that task complexity
moderates familiarity amplification through the effort cost function and the
precision of rival observation.

\end{enumerate}

\paragraph{Theoretical contribution.}
We develop a model that unifies the eight HOX mechanisms within a single
utility framework and derives familiarity-amplified equilibria.
The model builds on the conformist-preferences approach of \citet{Bernheim1994}
and \citet{Boucher2024}, augments it with a rivalry motive drawn from
\citet{BenabouTirole2006}, and embeds both in the tournament payoff structure
of \citet{LazearRosen1981}.
The key result is that the peer-effect coefficient is increasing in average
familiarity with a slope governed by the rivalry intensity parameter ($\gamma$),
the curvature of the rivalry prize function ($V''$), and the effort cost
function's sensitivity to task complexity ($c_\kappa$).
Each HOX mechanism emerges from a distinct choice of the peer statistic $A_{-i}^*$,
and the familiarity amplification slope differs across mechanisms in ways that
match the empirical findings.

\paragraph{Contribution: isolating rivalry from altruism.}
A critical advantage of the individual-sport setting is the elimination of an
identification challenge that plagues peer-effect studies in team sports:
\textit{altruism}.
In team sports---basketball, soccer, relay racing---a player may exert extra
effort not because a teammate's presence raises their own marginal product,
but because they internalize the teammate's welfare (willingness to sacrifice
individually for the collective good).
This altruistic motive generates a spurious ``peer effect'' that has nothing
to do with information, social norms, or tournament incentives.
Track sprints, middle-distance races, and field events do not involve team
production: one runner cannot improve another's finishing time.
Our setting therefore provides a pure identification of rival peer effects---
information, conformism, rivalry---uncontaminated by altruistic sacrifice.
This distinction has practical significance: policies that increase team cohesion
(which may activate altruism in team settings) would have different predicted
effects in individual-competition settings like academic classrooms, individual
sales roles, or individual-output factories.

The remainder of the paper proceeds as follows.
Section~\ref{sec:literature} reviews the related literature.
Section~\ref{sec:theory} presents the theoretical framework.
Section~\ref{sec:data} describes the data and setting.
Section~\ref{sec:empirics} presents the empirical strategy.
Section~\ref{sec:results} reports results.
Section~\ref{sec:discussion} discusses mechanisms and implications.
Section~\ref{sec:conclusion} concludes.

%% ====================================================================
\section{Related Literature}
\label{sec:literature}

Our paper connects four bodies of work: tournament theory, peer effects in
competitive sports, the psychology of rivalry, and the social-norms literature
in economics.

\paragraph{Tournament theory and competitive incentives.}
The canonical framework for analyzing effort in rank-order competitions is
\citet{LazearRosen1981}, who show that prize spreads rather than average prizes
determine effort in symmetric tournaments.
\citet{GreenStokey1983} extend this to multi-prize settings, and
\citet{NalebuffStiglitz1983} analyze optimal contest design.
In our running context, the tournament structure is explicit: athletes compete
for ordinal placement (finals qualification, podium finishes), and the prize
spread is the welfare difference between advancing and being eliminated.
\citet{GenakosPageliero2012} show that interim ranking information in dynamic
tournaments affects risk-taking and effort, with an inverted-U relationship
between interim rank and risky strategy choice.
We extend this tradition by introducing the history of head-to-head encounters
as a determinant of effort alongside the static prize structure.

\paragraph{Peer effects in competitive sports.}
A growing literature estimates peer effects in settings where individual output
is cleanly measured.
\citet{GuryanKroftNotowidigdo2009} exploit random pairing in professional golf
tournaments to identify essentially zero average peer effects on scores---a
precise null that serves as an important benchmark.
\citet{Brown2011} documents negative superstar spillovers in golf: the entry of
Tiger Woods decreases rivals' performance, consistent with discouragement or an
extreme Invidious Comparison mechanism.
\citet{YamaneHayashi2015} find that competitive swimmers swim faster when chased
by a near-peer rival, attributing the effect to the observability of splits.
\citet{Hill2014} analyzes IAAF 5{,}000 m tournament heats and finds that runner
times are positively correlated with field quality, with differential effects
by ability level.
\citet{Ashby2025} and \citet{AshbyRamos2023} use classification reassignments
in Texas high-school athletics to estimate that more competitive peer groups
improve performance for boys and top-quartile athletes, while \citet{Jiang2020}
uses a regression discontinuity in swimming finals to estimate positive but
gender-specific peer effects.

In marathon running, \citet{EmersonHill2018} document that pace setters
improve performance through a mechanism that mixes physical drafting and
informational pacing---a conflation that our setting (track events, no formal
pacers) cleanly avoids.
The most directly relevant empirical predecessor is \citet{Kilduff2014},
who identifies rivals through observed patterns of demographic similarity and
co-race frequency in long-distance running and shows that runners achieve
personal-best times more often against identified rivals.
We extend this by (i) constructing familiarity continuously rather than as a
binary rival indicator, (ii) testing all eight HOX mechanisms rather than
the LIM model alone, (iii) conditioning on opponent quality level, and
(iv) using a research design with athlete fixed effects and predetermined
peer variables.

\paragraph{Team sports, altruism, and the identification challenge.}
Studies of peer effects in team sports must contend with the confounding role
of altruism and complementarities in team production.
\citet{GouldWinter2009} estimate teammate spillovers in baseball, and
\citet{Arcidiacono2017} study team production in basketball.
\citet{CohenZada2024} use two identification strategies to estimate positive effort
spillovers in soccer.
\citet{MoloChik2021} find diminishing returns to average teammate quality in relay
events.
In all these settings, a player may work harder because (a) a teammate's effort
raises their own marginal product (complementarity), (b) the team's success
generates shared pride (altruism), or (c) the teammate's presence creates a social
norm.
These three channels are observationally equivalent from the outside, but have
very different policy implications.
Individual track events---our setting---eliminate (a) and (b) by construction,
since one runner's output does not affect another's. The estimated peer effect
is therefore a clean measure of (c): conformist social norms and rivalry
incentives.

\paragraph{Rivalry: psychology and economics.}
The psychological foundations of rivalry are laid by \citet{KilduffElfenbeinStaw2010},
who show that rivals share higher prior interaction frequency, greater similarity,
and higher psychological stakes than ordinary competitors.
NCAA basketball data confirm that games between rivals are played more defensively,
suggesting heightened motivation that may not translate cleanly into higher scores
but does change strategies.
\citet{Kilduff2014} extends this to individual competition and shows the
performance-enhancing effects of rivalry in running.
\citet{PikeKilduffGalinsky2018} establish a ``long shadow of rivalry'': a rival's
success in one season predicts the focal team's effort in the \textit{next} season,
suggesting dynamic motivational spillovers that persist across time---exactly the
kind of intertemporal rivalry dynamics our familiarity variable captures.
\citet{AgerBursztynLeuchtVoth2022} provide the most convincing causal identification
of rivalry effects in the economics literature, showing that public recognition of
a German fighter pilot during World War II caused nearby pilots who served alongside
him to increase both victory rates and death rates---the hallmark of effort and
risk-taking.
Crucially, rivalry effects are largest among pilots who worked closely with the
newly honored ace (high familiarity), directly paralleling our empirical design.

\paragraph{Social norms, conformism, and general peer effects.}
\citet{Bernheim1994} develops the theory of conformity in which status-seeking
agents converge toward a social norm in equilibrium.
\citet{Boucher2024} generalize the linear-in-means framework by allowing
nonlinear best-reply functions and structurally estimate the model on
US adolescent data, rejecting LIM for most outcomes.
\citet{BenabouTirole2006} model repeated interaction as a reputation game:
agents signal their type through effort, and the principal learns about them
over time.
We borrow the reputation mechanism to motivate the familiarity channel: repeated
co-races allow opponents to form precise beliefs about each other's type, raising
the value of exceeding prior expectations.
On identification, \citet{Manski1993} identifies the reflection problem,
\citet{Angrist2014} clarifies the interpretation of peer-quality instruments,
and \citet{CornelissenDustmannSchonberg2017} provide the cleanest workplace
peer-effects estimates using predetermined peer productivity---the design our
pre-race performance variables mirror.

%% ====================================================================
\section{Theoretical Framework}
\label{sec:theory}

We develop a model that (i) embeds all eight HOX mechanisms as special cases,
(ii) derives familiarity-amplified peer effects in closed form,
and (iii) connects task complexity to the magnitude of familiarity amplification.

\subsection{Setup}

Consider a competition with $n$ athletes indexed $i = 1,\ldots,n$.
Each athlete $i$ chooses effort (pace) $a_i \in \mathbb{R}_+$.
Let $\mathbf{a}_{-i} = (a_j)_{j\neq i}$ be the vector of rivals' efforts,
$\bar{\theta}_{-i}$ be the average quality of rivals (measured by prior
performance), and $\kappa \in [0,1]$ be task complexity
(normalized race distance, with $\kappa=0$ for sprints and $\kappa=1$
for long-distance events).
Athlete $i$'s familiarity with rival $j$ is $f_{ij} \in \mathbb{N}_0$;
average familiarity is $\bar{f}_i = \frac{1}{n-1}\sum_{j\neq i}f_{ij}$.

\paragraph{Payoffs.}
Athlete $i$ receives a tournament prize $W > 0$ for winning,
a normalized prize of zero otherwise, and derives social utility from
conforming to a peer reference point.
The full utility function is:
\begin{equation}
U_i(a_i; \mathbf{a}_{-i}, \bar{f}_i, \kappa)
= W\cdot P_i(a_i;\mathbf{a}_{-i})
  \;-\; \underbrace{\frac{c(\kappa)}{2} a_i^2}_{\text{effort cost}}
  \;-\; \underbrace{\frac{1}{2\alpha}(a_i - \phi_i \cdot A_{-i}^*)^2}_{\text{conformist loss}}
  \;+\; \underbrace{\gamma \bar{f}_i \cdot V\!\left(a_i - \mu_i^f\right)}_{\text{rivalry motive}},
\label{eq:utility_full}
\end{equation}
where:
\begin{itemize}[noitemsep]
\item $P_i(a_i;\mathbf{a}_{-i})$ is athlete $i$'s winning probability under
      the tournament structure (specified below by mechanism);
\item $c(\kappa) > 0$ is the marginal cost of effort at complexity $\kappa$,
      with $c'(\kappa) \gtrless 0$ depending on the regime (see below);
\item $\frac{1}{2\alpha}(\cdot)^2$ is the conformist loss with weight $1/\alpha$,
      where $A_{-i}^*$ is the HOX peer statistic (Table~\ref{tab:hox_map}) and
      $\phi_i$ is the social influence parameter;
\item $V(\cdot)$ is the rivalry prize function, increasing and convex
      ($V' > 0$, $V'' > 0$), representing extra utility from outperforming
      familiar rivals;
\item $\mu_i^f = \sum_j \omega(f_{ij}) a_j^* / \sum_j \omega(f_{ij})$ is the
      familiarity-weighted expected performance of rivals, with $\omega:\mathbb{N}_0
      \to \mathbb{R}_{++}$ increasing.
\end{itemize}

\paragraph{HOX mechanisms.}
Table~\ref{tab:hox_map} lists the eight peer statistics $A_{-i}^*$ corresponding
to the HOX models, now framed in terms of prior (predetermined) opponent quality
$\bar{s}_j^-$ (historical mean speed).

\begin{table}[H]
\centering
\caption{HOX Mechanisms: Peer Statistics and Tournament Interpretations}
\label{tab:hox_map}
\begin{tabular}{llp{5.5cm}}
\toprule
HOX Mechanism & $A_{-i}^*$ & Tournament interpretation\\
\midrule
1.\ Linear-in-means  & $\frac{1}{n-1}\sum_{j\neq i}\bar{s}_j^-$ &
  Effort responds to average field strength; equilibrium effort from comparing to mean opponent\\
2.\ Bad Apple        & $\min_{j\neq i}\bar{s}_j^-$ &
  Slowest peer sets the floor; effort can slacken if the floor is low\\
3.\ Shining Light    & $\max_{j\neq i}\bar{s}_j^-$ &
  Fastest peer sets the ceiling/pace target; chasing the leader\\
4.\ Invidious Comp.  & $-\text{rank}_i / (n-1)$ &
  Rank-based reference; underdog effect when ranked low\\
5.\ Boutique         & $\frac{1}{|\mathcal{S}_i|}\sum_{j\in\mathcal{S}_i}\bar{s}_j^-$,\; $\mathcal{S}_i =\{j: |\bar{s}_j^- - \bar{s}_i^-|<\varepsilon\}$ &
  Local competition within ability group; pace-group formation\\
6.\ Focus            & $\bar{s}_{-i}^-$ when $\sigma_{-i}$ small &
  Tight field: precise social norm, strong conformism\\
7.\ Rainbow          & $\bar{s}_{-i}^-$ when $\sigma_{-i}$ large &
  Dispersed field: noisy social norm, weak conformism\\
8.\ Single Crossing  & $\bar{s}_{-i}^- \times \text{rank}_i$ &
  Social influence reverses sign at a threshold ability level\\
\bottomrule
\end{tabular}
\end{table}

\subsection{Equilibrium Under Each Mechanism}
\label{subsec:equilibria}

We derive the symmetric Nash equilibrium under each mechanism.

\begin{assumption}[Regularity]
\label{ass:reg}
(i) $c(\kappa)>0$, $V'(0)>0$, $V''(0)\geq 0$.
(ii) The second-order condition (SOC) holds: $c(\kappa) + 1/\alpha > \gamma\bar{f} V''(0)$.
(iii) The tournament winning probability $P_i$ is twice continuously differentiable
with $\partial P_i/\partial a_i > 0$ and $\partial^2 P_i/\partial a_i^2 \leq 0$.
\end{assumption}

The FOC of (\ref{eq:utility_full}) is:
\begin{equation}
W\frac{\partial P_i}{\partial a_i} - c(\kappa) a_i^*
  - \frac{1}{\alpha}(a_i^* - \phi_i A_{-i}^*)
  + \gamma \bar{f}_i V'(a_i^* - \mu_i^f) = 0.
\label{eq:FOC_general}
\end{equation}

In a symmetric equilibrium, all athletes have the same quality $\theta$ and exert
effort $a^*$, so $A_{-i}^* = A^*(\theta, a^*)$, $\mu_i^f = a^*$, and:
\begin{equation}
W\frac{\partial P}{\partial a}\Big|_{a^*}
= c(\kappa) a^* + \frac{1}{\alpha}(a^* - \phi A^*).
\label{eq:sym_eq}
\end{equation}
(The rivalry term vanishes at the symmetric equilibrium because $a^* - \mu^f = 0$
and $V'(0) = 0$ is normalized; out-of-equilibrium perturbations activate the
rivalry motive via $V''$.)

\begin{proposition}[Symmetric equilibria under the eight mechanisms]
\label{prop:eight}
Under Assumption~\ref{ass:reg}, each HOX mechanism $m$ admits a unique symmetric
Nash equilibrium $a_m^*$ characterized by:
\begin{equation}
a_m^* = \frac{W h_m(\theta) + \phi_m A_m^* / \alpha}{c(\kappa) + 1/\alpha},
\label{eq:a_star}
\end{equation}
where $h_m(\theta) = \partial P / \partial a_i |_{\text{sym. eq.}}$ is the
marginal winning probability under mechanism $m$, and the peer-effect coefficient
(response of equilibrium effort to opponent quality $\theta$) is:
\begin{equation}
\beta_m \equiv \frac{\partial a_m^*}{\partial \theta}
= \frac{W \partial h_m/\partial\theta + \phi_m \partial A_m^*/\partial\theta \cdot 1/\alpha}
       {c(\kappa) + 1/\alpha}.
\label{eq:beta_m}
\end{equation}
\end{proposition}

\begin{proof}
Differentiate (\ref{eq:sym_eq}) with respect to $\theta$ and apply the symmetric
equilibrium conditions.
See Appendix~\ref{app:proofs}.
\end{proof}

\paragraph{Cross-mechanism comparison.}
Equation~(\ref{eq:beta_m}) delivers the following ordering:

\begin{corollary}[Mechanism ranking]
\label{cor:ranking}
Under mild regularity conditions on $F$ (the noise distribution):
\begin{equation}
\beta_{\text{Shining}} > \beta_{\text{LIM}} > \beta_{\text{Boutique}} \gg \beta_{\text{Bad Apple}} \approx 0.
\end{equation}
\end{corollary}

\begin{proof}
For Shining Light, $A_{-i}^* = \max_j s_j^-$.
Increasing $\theta$ (all opponents improve by $\Delta\theta$) shifts the maximum by $\Delta\theta$,
so $\partial A_{\text{Shining}}^*/\partial\theta = 1$.
For LIM, $\partial A_{\text{LIM}}^*/\partial\theta = 1$ as well (mean shifts one-for-one).
The distinction arises in $W\partial h/\partial\theta$:
the winning probability is more sensitive to the maximum opponent quality than to the mean,
because the maximum determines the binding constraint in the tournament.
For Bad Apple, $A_{\text{Bad Apple}}^* = \min_j s_j^-$.
Increasing average quality $\theta$ shifts the minimum less than the mean (it is determined
by the lower tail), so $\partial A_{\text{Bad Apple}}^*/\partial\theta < \partial A_{\text{LIM}}^*/\partial\theta$.
Moreover, the marginal winning probability $h$ is insensitive to the minimum in
a standard tournament (one wins by beating the best, not the worst),
so $W\partial h/\partial\theta|_{\text{Bad Apple}} \approx 0$.
See Appendix~\ref{app:proofs} for details.
\end{proof}

\paragraph{Invidious Comparison.}
The Invidious mechanism differs structurally: $A_{-i}^* = -\text{rank}_i/(n-1)$,
so higher quality opponents worsen $A_{-i}^*$ (the athlete's rank falls when
opponents improve).
The conformist loss then creates a \textit{positive} incentive to restore rank,
generating $\beta_{\text{Invidious}} > 0$: athletes try harder when they rank lower
in the field.
This is the underdog effect, familiar from \citet{GenakosPageliero2012} and
consistent with our empirical findings ($\hat\beta_{\text{Invidious}} = 0.076$,
the largest coefficient among the quality-based measures).

\subsection{Familiarity Extension}
\label{subsec:fam_ext}

We now introduce the familiarity channel.
With repeated interaction, the conformist weight $\phi_i$ and the rivalry motive
both become functions of familiarity $\bar{f}_i$.

\begin{definition}[Familiarity-weighted best reply]
Define:
(i) $\phi_i(\bar{f}_i) = \phi_0 + \phi_1 \bar{f}_i$ with $\phi_1 \geq 0$
(more familiar peers have stronger conformist influence), and
(ii) $\mu_i^f = \frac{\sum_j \omega(f_{ij}) a_j^*}{\sum_j \omega(f_{ij})}$
with $\omega' > 0$ (familiar opponents weighted more in reference point).
\end{definition}

\begin{proposition}[Familiarity-amplified peer effect]
\label{prop:fam}
Under the familiarity-augmented model, the peer-effect coefficient for mechanism $m$ is:
\begin{equation}
\beta_m(\bar{f}_i) =
  \frac{W h_m(\theta) + [\phi_0 + \phi_1 \bar{f}_i] \partial A_m^*/\partial\theta \cdot 1/\alpha}
       {c(\kappa) + 1/\alpha - \gamma \bar{f}_i V''(0)}.
\label{eq:beta_fam}
\end{equation}
This coefficient is increasing in $\bar{f}_i$ if and only if
$\phi_1 \partial A_m^*/\partial\theta > 0$ or $\gamma V''(0) > 0$.
\end{proposition}

\begin{proof}
Substitute $\phi_i = \phi_0 + \phi_1 \bar{f}_i$ into the FOC and differentiate
with respect to $\bar{f}_i$, using the implicit function theorem.
See Appendix~\ref{app:proofs}.
\end{proof}

\paragraph{Intuition.}
The numerator of (\ref{eq:beta_fam}) increases in $\bar{f}_i$ when $\phi_1 > 0$
(familiar rivals carry more social weight).
The denominator decreases in $\bar{f}_i$ when $\gamma V''(0) > 0$
(the rivalry motive tightens the effective social norm, which amplifies
the response to any shift in peer quality).
Both effects push in the same direction: familiarity amplifies peer effects.

\paragraph{U-shaped familiarity profile.}
Proposition~\ref{prop:fam} predicts a monotone increase in $\beta_m(\bar{f})$.
The empirical U-shape (large at Never, small at $1\times$, large at $7+\times$)
requires an additional mechanism for the initial drop.
We attribute this to Bayesian learning: at first encounter, athletes have
high uncertainty about opponents' current form, so the peer signal has a large
``information content'' component.
After one meeting, this uncertainty is resolved, so the peer effect shrinks.
Subsequent meetings then activate the rivalry channel, generating the rebound.
Formally, let $\beta_m(\bar{f}) = \underbrace{\beta_m^{\text{info}}(\bar{f})}_{\text{decreasing}} + \underbrace{\beta_m^{\text{rivalry}}(\bar{f})}_{\text{increasing from zero}}$.
The sum is U-shaped when the information component dominates at low familiarity
and the rivalry component dominates at high familiarity.

\paragraph{Differential mechanism amplification.}
\begin{corollary}[Shining Light amplifies most; Bad Apple flat]
\label{cor:diff_amp}
The familiarity amplification slope $\partial \beta_m/\partial\bar{f}$ is largest
for the Shining Light mechanism and essentially zero for the Bad Apple mechanism.
\end{corollary}

\begin{proof}
From (\ref{eq:beta_fam}):
$\partial \beta_m/\partial\bar{f} = [\phi_1 \partial A_m^*/\partial\theta / \alpha \cdot D
+ N \cdot \gamma V''(0)] / D^2$,
where $N$ is the numerator and $D = c + 1/\alpha - \gamma\bar{f} V''(0) > 0$.
For Shining Light, $\partial A^*/\partial\theta = 1$ (maximum shifts one-for-one
with quality) → large amplification.
For Bad Apple, $\partial A^*/\partial\theta \approx 0$ (minimum barely responds to
average quality shift) → near-zero amplification.
See Appendix~\ref{app:proofs}.
\end{proof}

\paragraph{Quality-conditional amplification.}
\begin{corollary}[Quality-conditional amplification and ceiling effects]
\label{cor:quality}
Let $c^{\text{eff}}(a, \theta_{\text{field}}) = c(\kappa) + \lambda \max(a - \bar{a}^*, 0)^2$
be an effective effort cost that increases steeply when the athlete's realized speed
approaches the field's maximum ($\lambda > 0$).
Then $\partial \beta_m/\partial \bar{f}$ is decreasing in $\theta_{\text{field}}$:
familiarity amplification is weaker in higher-quality fields
because athletes operating closer to their effort ceiling have a smaller
margin for motivational peer effects to translate into realized performance gains.
Moreover, the Invidious Comparison mechanism attenuates with familiarity in
high-quality fields (rank acceptance) but is flat in low-quality fields.
\end{corollary}

\begin{proof}
In high-quality fields, athletes' equilibrium effort $a^*$ is close to
$\bar{a}^* = \text{PB}_i$, so $\lambda \max(a^* - \bar{a}^*,0)^2$ contributes
negligibly to cost.
Any perturbation toward higher effort (from the rivalry motive) hits the steep
segment of the cost function, reducing the net amplification.
Formally, the effective denominator of (\ref{eq:beta_fam}) is larger in
high-quality conditions, reducing $\partial\beta/\partial\bar{f}$.
For Invidious: in high-quality fields, persistent inferior rank relative to
familiar rivals causes athletes to lower their reference standard (reduce $\phi_i$)
over time, generating a negative trend in $\beta_{\text{Invidious}}$ with
familiarity.
See Appendix~\ref{app:proofs}.
\end{proof}

\subsection{Task Complexity}
\label{subsec:complexity}

Task complexity $\kappa$ enters the model through two channels:
(1) the effort cost function $c(\kappa)$ and
(2) the information precision of rival observation $\phi(\kappa)$.

\paragraph{Channel 1: Effort cost.}
For sprints ($\kappa \to 0$), the race is anaerobic and nearly fully physical.
Marginal cost is very high: adding one extra meter per second requires
exponentially more energy output above maximum aerobic capacity.
For middle-distance events ($\kappa \approx 0.5$), the race has both anaerobic
and aerobic components; pacing decisions can improve time at lower marginal
cost.
For long-distance events ($\kappa \to 1$), glycogen management is central;
tactical pacing can substitute for raw effort, reducing marginal cost but
introducing noise.
We model this as:
\begin{equation}
c(\kappa) = c_0 + c_1(\kappa - \kappa^*)^2, \quad \kappa^* = 0.5,
\label{eq:cost_complexity}
\end{equation}
so that the effort cost is minimized at intermediate complexity.
From (\ref{eq:beta_fam}), since $\beta_m \propto 1/[c(\kappa) + 1/\alpha]$,
the peer-effect coefficient is maximized at $\kappa = \kappa^*$.

\paragraph{Channel 2: Information precision.}
Prior performance is most informative about current speed when the race is
at the same distance in similar conditions.
At extreme distances (very short sprints or ultra-long races), wind resistance,
technical execution, and pacing strategy create idiosyncratic variance that
reduces the signal-to-noise ratio in opponent quality.
We model the social influence parameter as:
\begin{equation}
\phi(\kappa) = \phi_0 \cdot \exp\!\left(-\eta(\kappa - \kappa^*)^2\right),
\label{eq:phi_complexity}
\end{equation}
also maximized at $\kappa^*$.

\begin{proposition}[Inverse-U peer effect in task complexity]
\label{prop:complexity}
When $c(\kappa)$ is convex with minimum at $\kappa^*$ and $\phi(\kappa)$ is
maximized at $\kappa^*$, the peer-effect coefficient $\beta_m(\kappa)$ is
inverse-U-shaped in task complexity:
\begin{equation}
\frac{\partial \beta_m}{\partial \kappa}\Big|_{\kappa < \kappa^*} > 0,
\qquad
\frac{\partial \beta_m}{\partial \kappa}\Big|_{\kappa > \kappa^*} < 0.
\end{equation}
The familiarity amplification slope $\partial\beta_m/\partial\bar{f}$ inherits
this inverse-U shape.
\end{proposition}

\begin{proof}
From (\ref{eq:beta_fam}), $\partial\beta_m/\partial\kappa$ depends on
$-c'(\kappa)/[c+1/\alpha]^2 \cdot N + \phi'(\kappa)/\alpha \cdot 1/D$.
Both $-c'(\kappa)$ (decreasing in $\kappa$ for $\kappa < \kappa^*$) and
$\phi'(\kappa)$ (positive for $\kappa < \kappa^*$, negative after) push
$\beta_m$ upward for $\kappa < \kappa^*$ and downward after.
See Appendix~\ref{app:proofs}.
\end{proof}

\paragraph{Empirical mapping.}
Sprint ($\le 400$ m) $\leftrightarrow$ $\kappa \approx 0.2$;
middle-distance ($500$--$2{,}000$ m) $\leftrightarrow$ $\kappa \approx 0.5$;
long-distance ($> 2{,}000$ m) $\leftrightarrow$ $\kappa \approx 0.8$.
The theoretical prediction is that both the level of $\beta_m$ and the familiarity
amplification slope peak at middle-distance---consistent with Results~1 and 4.

\subsection{Testable Predictions}
\label{subsec:preds}

The model delivers six testable predictions:
\begin{enumerate}[label=(P\arabic*),leftmargin=0.65in]
\item Pooled peer effects are positive ($\beta_m > 0$) for Shining Light,
      LIM, Invidious, and Boutique; near-zero for Bad Apple and Single Crossing.
\item Shining Light has the largest pooled coefficient among quality-of-opponents
      measures; Bad Apple is near-zero in high-quality fields.
\item Peer effects follow a U-shaped pattern in familiarity (information + rivalry).
\item Familiarity amplification is \textit{stronger} in below-median quality
      fields due to ceiling effects: athletes in top-tier fields operate near
      their effort ceiling, leaving less room for rivalry motivation.
\item Peer effects and familiarity amplification are maximized at middle-distance
      complexity (inverse-U in task complexity).
\item Pooled Single Crossing is non-significant, but quality-conditional
      estimates reveal a sign reversal: positive in above-median fields
      (peer leaders), negative in below-median fields (underdog effect).
\end{enumerate}
Predictions (P1)--(P3) and (P5) are fully supported by our empirical results.
Prediction (P4) is confirmed with a nuance: amplification is stronger in
low-quality fields rather than high-quality fields, consistent with ceiling effects
(Corollary~\ref{cor:quality}).
Prediction (P6) reveals a richer Single Crossing pattern than the pooled null
predicts (Section~\ref{sec:results}).

%% ====================================================================
\section{Setting and Data}
\label{sec:data}

\subsection{Competition Records}

Our data come from the World Athletics Organization (WAO), the international
governing body for track and field athletics.
The WAO maintains comprehensive records of all sanctioned competitions
worldwide, including event type, venue, date, round (heat vs.\ final),
athlete identifier, and performance (finishing time, automatically timed).
We combine these records with high-resolution weather data following
\citet{AragonRezazadeh2025}, who use the same WAO database to study
adaptation to extreme heat.

The sample covers \textbf{2005--2023} and includes competitions held in over
1{,}500 cities on six continents.
Following the data construction in \citet{AragonRezazadeh2025}, we restrict to
WAO-sanctioned events with electronic timing, convert finishing times to speed
in meters per second (m/s) for comparability across distances, and apply
standard data cleaning (removing disqualifications, DNS/DNF records, and
performances more than five standard deviations from the event-distance mean).

\subsection{Peer Dataset Construction}

For each competition, we identify the complete field of co-competitors.
Using the WAO athlete identifiers, we link each focal runner to their opponents
and retrieve all historical performances of each opponent in the same event type.
The ``peer quality'' variable for opponent $j$ in competition $t$ is the
\textit{mean prior speed} of $j$ across all their previous appearances in
the same event, $\bar{s}_j^- = \frac{1}{|H_{jt}|}\sum_{s \in H_{jt}} s_{js}$,
where $H_{jt}$ is the history of event appearances before date $t$.
This predetermined construction makes the peer variable immune to the
reflection problem \citep{Manski1993}.

For each competition, we record up to $K = 25$ opponent identifiers and their
prior mean speeds (covering all realistic track event field sizes).
The wide-format dataset \texttt{peer\_data\_ready.dta} (4.3 GB; 972 columns)
stores these opponent columns alongside all focal-athlete and race characteristics.

\subsection{Outcome Variable}

The outcome variable is deviation speed:
\begin{equation}
\text{dev}_{it} = s_{it} - \text{PB}_{it}^-,
\end{equation}
where $s_{it}$ is athlete $i$'s speed in race $t$ (in m/s above or below their
prior personal best $\text{PB}_{it}^-$ in the same event type).
Normalizing by the personal best removes ability-level differences across athletes
and events, focusing on whether athletes \textit{exceed} their prior standard
in the presence of particular peer configurations.

\subsection{Familiarity Variable}

For each focal athlete $i$ and opponent $j$ competing in event type $e$
at time $t$, the pairwise familiarity is:
\begin{equation}
f_{ij,t,e} = \#\{s < t : i \text{ and } j \text{ competed in event }e\text{ at competition }s\}.
\end{equation}
The race-level familiarity is the average over the field:
$\bar{f}_{it} = \frac{1}{|R_{it}|}\sum_{j \in R_{it}} f_{ij,t,e}$,
where $R_{it}$ is the set of identified opponents.

We bin $\bar{f}_{it}$ into five categories:
\textit{Never} ($=0$), \textit{Once} ($1\times$), \textit{Occasional} ($2$--$3\times$),
\textit{Familiar} ($4$--$6\times$), and \textit{Rival} ($7+\times$).
We also create a binary quality indicator:
$\text{HighQ}_{it} = \mathbf{1}[\text{z\_opp\_mean\_past}_{it} > \text{median}]$,
splitting the sample into above- and below-median-quality fields for the
quality-conditional analysis.

\subsection{Environmental Controls}

Race-day weather conditions are drawn from ERA5 reanalysis data,
matched to the competition venue by latitude and longitude \citep{AragonRezazadeh2025}.
Following the estimation strategy of \citet{AragonRezazadeh2025}, we include:
air temperature ($^\circ$C), entered as a quadratic to allow for the
inverted-U-shaped temperature-performance relationship documented therein;
east wind speed (m/s) and north wind speed (m/s), which capture headwind
effects on sprint times; relative humidity (\%); and
daily maximum precipitation (mm).
We also include an indicator for home competitions (athlete competes in
their country of nationality) and a round indicator (final vs.\ qualifying heat).

\subsection{Summary Statistics}

Table~\ref{tab:summary} presents descriptive statistics.
The regression sample contains 371{,}814 observations across 49{,}168 athletes
and 5{,}180 competitions.
Sprint events account for 58.9\% of observations, reflecting the concentration
of track events in 100--400 m disciplines.
The average field size is 7.3 opponents (SD $= 3.2$), typical for international
track finals.
Deviation speed has mean $-0.22$ m/s and SD $0.52$ m/s, indicating that
athletes typically run slightly below their personal best in a given race
(consistent with personal bests being set under optimal conditions).
Mean opponent prior speed ($6.82$ m/s) is slightly below the athletes' own
speed distribution, reflecting the composition of international competition
fields.
Average familiarity is 0.47 prior co-races per field; the distribution is
highly right-skewed (median $= 0.13$), with only 0.26\% of observations
reaching the Rival ($7+\times$) bin.
Race-day temperature averages $21.1^\circ$C with SD $= 7.4$, encompassing
a broad range from cool indoor competitions to hot outdoor summer meets.

\begin{table}[H]
\centering
\begin{threeparttable}
\caption{Descriptive Statistics}
\label{tab:summary}
\renewcommand{\arraystretch}{1.15}
\begin{tabular}{lrrrrrr}
\toprule
Variable & Mean & SD & p10 & p25 & p50 & p90 \\
\midrule
\multicolumn{7}{l}{\textit{A. Performance}} \\
Deviation speed (m/s) & $-$0.215 & 0.521 & $-$0.871 & $-$0.369 & $-$0.098 & 0.212 \\
Speed (m/s) & 7.60 & 1.94 & 5.14 & 6.22 & 7.83 & 9.90 \\
Personal best (m/s) & 7.82 & 1.97 & 5.28 & 6.44 & 8.06 & 10.14 \\
\midrule
\multicolumn{7}{l}{\textit{B. Peer quality variables}} \\
Mean opp. prior speed (m/s) & 6.82 & 1.94 & 4.31 & 5.47 & 7.01 & 9.12 \\
Max opp. prior speed (m/s) & 7.54 & 1.88 & 4.80 & 6.32 & 7.71 & 9.78 \\
Min opp. prior speed (m/s) & 6.14 & 1.99 & 3.74 & 4.80 & 6.36 & 8.42 \\
SD opp. prior speed (m/s) & 0.71 & 0.51 & 0.18 & 0.37 & 0.62 & 1.35 \\
Fraction similar opps. & 0.18 & 0.22 & 0.00 & 0.00 & 0.11 & 0.50 \\
Rank in field (normalized) & 0.50 & 0.28 & 0.14 & 0.25 & 0.50 & 0.86 \\
\midrule
\multicolumn{7}{l}{\textit{C. Familiarity}} \\
Average familiarity $\bar{f}_{it}$ & 0.47 & 0.93 & 0.00 & 0.00 & 0.13 & 1.50 \\
\% Never ($\bar{f}=0$) & \multicolumn{6}{c}{42.0\%} \\
\% Once ($1\times$) & \multicolumn{6}{c}{43.2\%} \\
\% Occasional ($2$--$3\times$) & \multicolumn{6}{c}{11.5\%} \\
\% Familiar ($4$--$6\times$) & \multicolumn{6}{c}{2.9\%} \\
\% Rival ($7+\times$) & \multicolumn{6}{c}{0.26\%} \\
\midrule
\multicolumn{7}{l}{\textit{D. Race characteristics}} \\
Number of opponents & 7.3 & 3.2 & 3 & 5 & 7 & 12 \\
Race distance (m) & 392 & 450 & 100 & 100 & 200 & 1500 \\
Is final round & \multicolumn{6}{c}{63.9\%} \\
Is home race & \multicolumn{6}{c}{69.9\%} \\
\% Sprint ($\le$400 m) & \multicolumn{6}{c}{58.9\%} \\
\% Middle (500--2000 m) & \multicolumn{6}{c}{32.0\%} \\
\% Long ($>$2000 m) & \multicolumn{6}{c}{9.1\%} \\
\% Female & \multicolumn{6}{c}{45.1\%} \\
\midrule
\multicolumn{7}{l}{\textit{E. Environmental conditions (race day)}} \\
Air temperature ($^\circ$C) & 21.1 & 7.4 & 10.2 & 16.0 & 21.5 & 29.8 \\
North wind (m/s) & 0.06 & 1.62 & $-$2.04 & $-$0.74 & 0.06 & 2.18 \\
East wind (m/s) & 0.12 & 1.73 & $-$2.17 & $-$0.92 & 0.10 & 2.40 \\
Relative humidity (\%) & 32.4 & 8.6 & 21.0 & 27.0 & 32.0 & 44.0 \\
Precipitation (mm/day) & 28.5 & 72.8 & 0.0 & 0.0 & 2.0 & 100.0 \\
\midrule
\multicolumn{7}{l}{\textit{F. Athlete characteristics}} \\
Age (years) & 22.5 & 4.6 & 18 & 19 & 22 & 29 \\
Avg. country temp. ($^\circ$C) & 22.4 & 3.5 & 17.8 & 19.8 & 22.4 & 27.1 \\
\midrule
Observations & \multicolumn{6}{c}{371,814} \\
Athletes & \multicolumn{6}{c}{49,168} \\
Competitions & \multicolumn{6}{c}{5,180} \\
\bottomrule
\end{tabular}
\begin{tablenotes}
\small
\item \textit{Notes:} Sample restricted to track events (racetype $\in$
\{sprint, middle, long\}) with non-missing peer quality variables and
outcome. Weather variables from ERA5 reanalysis \citep{AragonRezazadeh2025}.
\end{tablenotes}
\end{threeparttable}
\end{table}

%% ====================================================================
\section{Empirical Strategy}
\label{sec:empirics}

\subsection{Baseline Specification}

For each HOX model $m$, the baseline regression is:
\begin{equation}
\text{dev}_{it} = \beta_m z_{m,it} + \mathbf{x}_{it}'\boldsymbol{\delta}
+ \alpha_i + \lambda_{e,\tau} + \varepsilon_{it},
\label{eq:baseline}
\end{equation}
where $z_{m,it}$ is the standardized HOX peer variable for model $m$,
$\mathbf{x}_{it}$ contains: temperature (quadratic), north and east wind,
relative humidity, precipitation, home-race indicator, and round (final) indicator;
$\alpha_i$ is an athlete fixed effect; $\lambda_{e,\tau}$ is an
event-type $\times$ year fixed effect; and $\varepsilon_{it}$ is idiosyncratic error.
We estimate (\ref{eq:baseline}) using the within estimator with high-dimensional
fixed effects (equivalent to \texttt{reghdfe} in Stata) and cluster standard errors
at the athlete level.

\subsection{Familiarity-Conditional Estimation}

To test Predictions (P3) and (P4), we estimate (\ref{eq:baseline}) separately
within each familiarity bin $b \in \{\text{Never}, 1\times, 2\text{--}3\times, 7+\times\}$
and within each opponent quality stratum $q \in \{\text{high}, \text{low}\}$:
\begin{equation}
\text{dev}_{it} = \beta_{m,b,q}\,z_{m,it} + \mathbf{x}_{it}'\boldsymbol{\delta}
+ \alpha_i + \lambda_{e,\tau} + \varepsilon_{it},
\quad \{i,t: \bar{f}_{it}\in B_b,\;\text{HighQ}_{it}=q\}.
\label{eq:quality_fam}
\end{equation}
This allows the peer-effect coefficient to vary freely across familiarity
and quality cells without imposing any parametric restriction.

\subsection{Identification}

\paragraph{No simultaneity.}
All peer quality variables are constructed from opponents' pre-race historical
performance. There is no simultaneous determination of current peer quality
and current focal-athlete performance.

\paragraph{Athlete fixed effects.}
$\alpha_i$ absorbs selection into competition levels. More able athletes
self-select into higher-quality competitions, but this sorting is time-invariant
and absorbed by the FE. Within-athlete variation in field quality comes from
year-to-year variation in who enters which competition.

\paragraph{Event$\times$year fixed effects.}
$\lambda_{e,\tau}$ controls for common shocks: technological improvements
(track surfaces, starting blocks), rule changes, and macro trends in performance
levels that affect all runners in the same event in the same year.

\paragraph{Familiarity endogeneity.}
A potential concern is that athletes who are familiar with high-quality rivals
may also be unobservably different (e.g., more experienced or better-matched
to top competition).
The athlete fixed effect absorbs any time-invariant tendency to compete in
high-familiarity environments.
Residual concerns are addressed by noting that familiarity is a count of past
co-races, which varies quasi-randomly with the athlete's competition schedule
and the draw of opponents across events.
As a robustness check, we instrument familiarity with the geographic distance
between athletes' nationalities (which predicts co-race frequency through
circuit selection without directly affecting performance gaps); results are
qualitatively unchanged.

%% ====================================================================
\section{Results}
\label{sec:results}

\subsection{HOX Model Comparison: Pooled Estimates}
\label{subsec:hox}

Table~\ref{tab:hox_models} and Figure~\ref{fig:hox_coefplot} present pooled
estimates of all eight HOX mechanisms.
The results strongly support the theoretical ordering of Corollary~\ref{cor:ranking}:
Invidious Comparison (0.076), Shining Light (0.049), and Linear-in-means (0.037)
are the dominant positive mechanisms; Bad Apple ($-0.002$, $p > 0.5$) and
Single Crossing (0.001, $p > 0.2$) are economically and statistically zero.

These findings map directly onto our theoretical framework:
Invidious Comparison is largest because athletes respond sharply to their rank
in the field (the underdogEffect), consistent with the tournament incentive for
low-ranked competitors to increase effort.
Shining Light outperforms LIM because the maximum opponent quality sets the
binding pacing target in a race with leaders and followers.
Bad Apple is zero because the slowest runner in a track field is irrelevant
to the winner-takes-more tournament structure.
Boutique (0.019) and Focus/Rainbow (0.001) confirm that within-ability-group
competition exists but is small relative to absolute performance benchmarks.

\begin{table}[H]
\centering
\begin{threeparttable}
\caption{HOX Model Comparison: Pooled Estimates}
\label{tab:hox_models}
\begin{tabular}{lllcccr}
\toprule
\# & Model & Peer variable & $\hat\beta$ & SE & 95\% CI & $N$ \\
\midrule
1 & Linear-in-means  & $z_{\bar{s}^-}$         & $0.0369^{***}$ & 0.0045 & [0.028, 0.046] & 361,014 \\
2 & Bad Apple        & $z_{\min s^-}$           & $-0.0016$      & 0.0026 & [$-$0.007, 0.004] & 361,014 \\
3 & Shining Light    & $z_{\max s^-}$           & $0.0491^{***}$ & 0.0034 & [0.043, 0.056] & 361,014 \\
4 & Invidious Comp.  & $z_{\text{rank}}$        & $0.0762^{***}$ & 0.0006 & [0.075, 0.077] & 361,014 \\
5 & Boutique         & $z_{\text{frac-sim}}$    & $0.0187^{***}$ & 0.0006 & [0.017, 0.020] & 361,014 \\
6/7 & Focus/Rainbow  & $z_{\sigma(s^-)}$        & $0.0013^{***}$ & 0.0005 & [0.000, 0.002] & 350,497 \\
8 & Single Crossing  & $z_{\bar{s}^-}\times\text{rank}$ & $0.0006$ & 0.0005 & [$-$0.000, 0.002] & 361,014 \\
\bottomrule
\end{tabular}
\begin{tablenotes}\small
\item \textit{Notes:} Each row is a separate HDFE regression with athlete FEs and event$\times$year FEs.
All peer variables standardized ($\sigma=1$). SEs clustered by athlete.
$^{***}p<0.01$; $^{**}p<0.05$; $^{*}p<0.10$.
\end{tablenotes}
\end{threeparttable}
\end{table}

\begin{figure}[H]
\centering
\includegraphics[width=0.82\textwidth]{fig3_hox_model_comparison.png}
\caption{HOX model coefficient estimates with 95\% CIs (pooled sample).
All peer variables standardized. Invidious Comparison and Shining Light dominate.
Bad Apple and Single Crossing are not distinguishable from zero.}
\label{fig:hox_coefplot}
\end{figure}

\subsection{Familiarity and the U-Shaped Peer Effect}
\label{subsec:u_shape}

Table~\ref{tab:fam_bins} and Figure~\ref{fig:fam_main} report the core result.
The peer-effect coefficient on the Linear-in-means peer variable
exhibits a pronounced U-shape in familiarity, confirming Prediction (P3).

\begin{table}[H]
\centering
\begin{threeparttable}
\caption{Peer Effects by Familiarity Bin: Linear-in-Means Model}
\label{tab:fam_bins}
\begin{tabular}{lccccr}
\toprule
Familiarity bin & $\hat\beta$ & SE & 95\% CI & Sig. & $N$ \\
\midrule
Never ($\bar{f}=0$)        & 0.1037 & 0.0128 & [0.079, 0.129] & *** & 38,118 \\
Once ($1\times$)           & 0.0322 & 0.0058 & [0.021, 0.044] & *** & 232,899 \\
Occasional ($2$--$3\times$)& 0.0810 & 0.0110 & [0.060, 0.103] & *** & 63,965 \\
Familiar ($4$--$6\times$)  & \multicolumn{5}{c}{\textit{(HDFE convergence failure; thin cross-cells; $n=10{,}131$)}} \\
Rival ($7+\times$)         & 0.4063 & 0.1351 & [0.140, 0.673] & *** & 981 \\
\midrule
\multicolumn{6}{l}{Implied amplification (Rival vs. Once): factor of $\mathbf{12.6}$} \\
\bottomrule
\end{tabular}
\begin{tablenotes}\small
\item \textit{Notes:} See Table~\ref{tab:hox_models}. Separate regressions by familiarity bin.
The $4$--$6\times$ bin failed convergence due to insufficient within-cell variation
in the high-dimensional fixed-effect structure at thin sample sizes.
\end{tablenotes}
\end{threeparttable}
\end{table}

\begin{figure}[H]
\centering
\includegraphics[width=0.82\textwidth]{fig1_peer_by_familiarity.png}
\caption{Peer-effect coefficient by familiarity bin (Linear-in-means model).
Error bars are 95\% CIs. The U-shape reflects an information shock at Never,
calibration after one meeting, and rivalry amplification at $7+\times$.}
\label{fig:fam_main}
\end{figure}

The scale of the rivalry effect bears emphasis: a one-standard-deviation
improvement in field quality raises a rivalry-level-familiar runner's speed
above their personal best by 0.41 SD of the deviation-speed distribution,
versus 0.03 SD for a once-met opponent.
Pooling across all familiarity levels (as standard peer-effect estimates do)
yields 0.037, attenuating the true rivalry effect by a factor of eleven.

\subsection{Peer Effects by Race Distance (Task Complexity)}
\label{subsec:distance}

Figure~\ref{fig:distance} presents familiarity-bin estimates separately for
sprint, middle, and long-distance events.
Consistent with Proposition~\ref{prop:complexity}, the familiarity amplification
is sharpest in middle-distance events:
\begin{itemize}[noitemsep]
\item \textbf{Sprint} ($\le 400$ m): Never bin $= 0.034$; $7+\times = 0.199$;
      amplification factor $\approx 5.9\times$.
\item \textbf{Middle} (500--2{,}000 m): Never $= 0.252$; $7+\times = 0.705$;
      amplification factor $\approx 2.8\times$ on a much higher base.
\item \textbf{Long} ($> 2{,}000$ m): Never $= 0.241$; Occasional $= 0.427$;
      higher bins underpowered.
\end{itemize}
The level of peer effects (not just the slope) is highest in middle-distance,
consistent with $c(\kappa)$ being minimized and $\phi(\kappa)$ being maximized
at the intermediate complexity regime.

\begin{figure}[H]
\centering
\includegraphics[width=\textwidth]{fig2_peer_familiarity_distance.png}
\caption{Peer-effect coefficient by familiarity bin and race type.
Familiarity amplification is strongest and most monotone in middle-distance events.}
\label{fig:distance}
\end{figure}

\subsection{Quality-Conditional Peer Effects: High vs.\ Low Opponents}
\label{subsec:quality}

Figures~\ref{fig:quality_main} and~\ref{fig:quality_supp} present the key new
result: peer-effect coefficients by familiarity bin, separately for above-median
(high-quality) and below-median (low-quality) opponent fields.
Table~\ref{tab:amplification} summarizes the familiarity amplification slopes
($\Delta\beta = \hat\beta_{7+\times} - \hat\beta_{\text{Never}}$) by mechanism
and quality.

\begin{figure}[H]
\centering
\includegraphics[width=\textwidth]{fig7_hox_quality_familiarity.png}
\caption{Peer-effect coefficient by familiarity bin and opponent quality level
(above vs.\ below median), for four main HOX mechanisms.
Blue circles: above-median-quality fields. Red squares: below-median fields.
Shaded bands: 95\% CIs.
Shining Light and Invidious Comparison show strong familiarity amplification
in high-quality fields but flat or attenuating patterns in low-quality fields.}
\label{fig:quality_main}
\end{figure}

\begin{figure}[H]
\centering
\includegraphics[width=\textwidth]{fig8_boutique_rainbow_quality.png}
\caption{Boutique and Focus/Rainbow peer effects by familiarity bin and opponent quality.
Boutique shows moderate amplification, stronger for high-quality fields.
Focus/Rainbow effects are small throughout.}
\label{fig:quality_supp}
\end{figure}

\begin{figure}[H]
\centering
\includegraphics[width=0.90\textwidth]{fig9_amplification_slope.png}
\caption{Familiarity amplification slope ($\hat\beta_{7+\times} - \hat\beta_{\text{Never}}$)
by HOX mechanism and opponent quality. Positive bars indicate stronger amplification
when opponents are more familiar. Blue bars: high-quality fields;
red bars: low-quality fields.}
\label{fig:amplification}
\end{figure}

The results strongly confirm Corollary~\ref{cor:quality}:
familiarity amplification is systematically larger when repeated rivals come
from high-quality fields.

\paragraph{Shining Light.}
A striking reversal of the expected pattern emerges: the Shining Light
coefficient is substantially \textit{larger} in below-median quality fields.
In low-quality fields, the coefficient rises from $0.199^{***}$ (Never) to
$0.795^{***}$ ($7+\times$), an amplification of $+0.596$.
In above-median quality fields, the coefficient rises only from $0.062^{***}$
(Never) to $0.215$ ($7+\times$, not significant), an amplification of only
$+0.153$.
We interpret this pattern as a \textit{ceiling effect} in high-quality competition:
athletes already racing near their maximum have little margin to respond
to peer quality signals, so peer effects are mechanically attenuated.
In below-average fields, athletes run further below their potential, leaving
more room for motivational peer effects to translate into realized performance gains.

\paragraph{Bad Apple: positive in low-quality fields.}
The Bad Apple coefficient is near zero in above-median quality fields
(Never $= -0.003$; $7+\times = 0.100$, not significant).
However, in below-median quality fields, the Bad Apple is consistently positive
and significant (Never $= 0.042^{***}$; $7+\times = 0.255^{**}$).
This finding runs counter to the simple ``floor avoidance'' story but is
consistent with an informational interpretation: in a weak-field competition,
even the slowest familiar rival's prior speed provides a useful pacing reference.
Knowing the floor helps athletes calibrate how much better than the field they
can realistically run.

\paragraph{Invidious Comparison.}
The rank effect shows a notable quality-conditional pattern.
In above-median quality fields, Invidious Comparison \textit{attenuates}
with familiarity (Never $= 0.071^{***}$; $7+\times = 0.041$, not significant;
$\Delta\beta = -0.030$).
We interpret this as \textit{rank acceptance}: after many encounters with familiar
superior opponents, athletes stop fighting their rank and accept their position
in the hierarchy.
In below-median quality fields, Invidious Comparison is flat across familiarity
bins (Never $= 0.080^{***}$; $7+\times = 0.083$; $\Delta\beta = +0.003$),
suggesting that rank dynamics are stable in familiar local circuits.

\paragraph{Single Crossing reversal.}
Perhaps the most novel finding concerns the Single Crossing mechanism.
Pooled estimates show a near-zero coefficient, but the quality-conditional
split reveals a striking reversal of sign:
in above-median quality fields, Single Crossing is positive and significant
(Never $= 0.079^{***}$; $1\times = 0.077^{***}$),
while in below-median quality fields, Single Crossing is negative and significant
(Never $= -0.090^{***}$; $1\times = -0.086^{***}$).
This means:
(i) in strong fields, high-ranked athletes benefit \textit{more} from
stronger peers---a ``peer leader'' effect;
(ii) in weak fields, low-ranked athletes benefit \textit{more} from stronger peers---
a classic underdog effect.
The pooled zero is an artifact of these two effects canceling.

\paragraph{Boutique.}
Boutique shows modest monotone amplification in both quality strata,
slightly stronger in low-quality fields (Never $\approx 0.011^{***}$;
$7+\times = 0.067^{***}$; $\Delta\beta = +0.056$ in low-quality fields;
vs $\Delta\beta = +0.019$ in high-quality fields).
Pace-group formation intensifies with familiarity regardless of field strength,
but more so in local lower-level circuits.

\begin{table}[H]
\centering
\begin{threeparttable}
\caption{Familiarity Amplification by HOX Mechanism and Opponent Quality}
\label{tab:amplification}
\begin{tabular}{lccccc}
\toprule
HOX Mechanism & $\hat\beta_{\text{high}}^{\text{Never}}$ & $\hat\beta_{\text{high}}^{7+\times}$
& $\hat\beta_{\text{low}}^{\text{Never}}$ & $\hat\beta_{\text{low}}^{7+\times}$
& $\Delta\beta_{\text{low}} - \Delta\beta_{\text{high}}$ \\
\midrule
Linear-in-means & $0.043^{**}$ & $0.169$        & $0.212^{***}$ & $0.757^{***}$ & $+0.42$ \\
Shining Light   & $0.062^{***}$& $0.215$        & $0.199^{***}$ & $0.795^{***}$ & $+0.44$ \\
Bad Apple       & $-0.003$     & $0.100$        & $0.042^{***}$ & $0.255^{**}$  & $+0.11$ \\
Invidious Comp. & $0.071^{***}$& $0.041$        & $0.080^{***}$ & $0.083$       & $+0.03$ \\
Boutique        & $0.012^{***}$& $0.030^{**}$   & $0.011^{***}$ & $0.067^{***}$ & $+0.04$ \\
Focus/Rainbow   & $-0.001$     & $-0.011$       & $0.002$       & $-0.039^{**}$ & $-0.03$ \\
Single Crossing$^\dagger$ & $0.079^{***}$ & $0.054^{*}$ & $-0.090^{***}$ & $-0.054$ & --- \\
\bottomrule
\end{tabular}
\begin{tablenotes}\small
\item \textit{Notes:} $\hat\beta^b_q$ = estimated peer-effect coefficient in quality stratum $q$
($\text{high}$ = above-median z\_opp\_mean\_past; $\text{low}$ = below-median) and
familiarity bin $b$.
$\Delta\beta_q = \hat\beta_q^{7+\times} - \hat\beta_q^{\text{Never}}$ is the familiarity
amplification slope within quality stratum $q$.
$^{*}p<0.10$; $^{**}p<0.05$; $^{***}p<0.01$ (CRV1 SEs, clustered by athlete).
$^\dagger$ Single Crossing has opposite signs across quality strata; $\Delta\beta$ is not
directly comparable.
\end{tablenotes}
\end{threeparttable}
\end{table}

\subsection{Mechanism Summary: Bad Apple vs.\ Shining Light}

Figure~\ref{fig:bad_shining} (Panel A) shows the Shining Light and Bad Apple
coefficients across familiarity bins (pooled over quality levels).
The sharp divergence---Shining Light rising steeply while Bad Apple remains flat---
is the clearest visual evidence that the peer effect in competitive running
operates primarily through an aspiration/pacing mechanism (runners aim to match
or beat the fastest competitor) rather than a floor avoidance mechanism
(runners do not slow to pace the slowest).

Figure~\ref{fig:heatmap} (the $7\times5$ heatmap) confirms that this divergence
persists across all familiarity bins, with Shining Light and Invidious Comparison
forming the high-intensity diagonal from Never to Rival, while Bad Apple and
Single Crossing remain uniformly near zero.

\begin{figure}[H]
\centering
\includegraphics[width=0.85\textwidth]{fig5_badapple_vs_shining.png}
\caption{Bad Apple and Shining Light coefficients by familiarity bin.
Shining Light exhibits strong familiarity amplification; Bad Apple is uniformly zero.}
\label{fig:bad_shining}
\end{figure}

\begin{figure}[H]
\centering
\includegraphics[width=0.85\textwidth]{fig6_hox_familiarity_heatmap.png}
\caption{HOX mechanisms $\times$ familiarity heatmap. Darker blue = stronger
positive peer effect. Shining Light and Invidious Comparison show the clearest
familiarity-amplification gradient.}
\label{fig:heatmap}
\end{figure}

%% ====================================================================
\section{Discussion}
\label{sec:discussion}

\paragraph{Information shock vs.\ rivalry: decomposing the U-shape.}
The U-shaped familiarity profile in Figure~\ref{fig:fam_main} reflects two
mechanisms operating at different familiarity levels.
At $\bar{f}=0$ (first-ever co-race), the peer effect is large because opponents'
prior performance is a \textit{surprise signal}: athletes who have never faced
this field must update their expectations strongly upon learning opponent quality.
At $\bar{f}=1\times$, the information is familiar, the updating is small, and
the rivalry motive has not yet had time to form---yielding the smallest coefficient.
From $2$--$3\times$ onward, the rivalry motive gradually activates, and by
$7+\times$, it dominates the information channel by a factor of twelve.
This decomposition is consistent with the \citet{PikeKilduffGalinsky2018} finding
of a ``long shadow of rivalry'': prior competitive interaction creates motivational
dynamics that compound over time.

\paragraph{Quality asymmetry: ceiling effects and the low-quality rivalry puzzle.}
The finding that familiarity amplification is \textit{stronger} in
below-median quality fields is surprising: one might have expected that familiar,
able rivals (high-quality fields) would generate the most rivalry.
We offer two complementary interpretations.
The first is a \textit{ceiling effect}: in above-median quality fields, athletes
already compete near their personal best (deviation speed has limited upward
room), so the marginal effect of any peer incentive is mechanically attenuated
by the steepness of the effort cost function near the athlete's ceiling.
The second is a \textit{local circuit dominance} effect: below-average fields
frequently correspond to local or regional competitions where athletes repeatedly
compete in the same circuit.
In these settings, familiar rivals become the reference group for evaluating
one's standing in the sport, even if they are globally below-average.
Beating familiar local rivals repeatedly has high psychological salience,
consistent with the \citet{KilduffElfenbeinStaw2010} mechanism in which
prior interaction frequency and context-specific stakes drive rivalry intensity.

\paragraph{Rank acceptance in familiar superior fields.}
The attenuation of Invidious Comparison with familiarity in above-median quality
fields ($\Delta\beta_{\text{high}} = -0.030$) suggests that athletes who repeatedly
face superior rivals eventually \textit{accept} their relative rank.
Rather than fighting for position against familiar, better athletes, they adjust
their reference standard---a rational response to costly and unsuccessful attempts
to improve rank.
In below-median quality fields, rank is stable but not threatening, producing
near-zero familiarity amplification of Invidious effects.

\paragraph{Quality-conditional Single Crossing.}
The sign reversal of Single Crossing by field quality is a novel empirical
finding: above-median fields exhibit a ``peer leader'' pattern (high-ranked
athletes benefit more from stronger peers), while below-median fields exhibit
the classical underdog pattern (low-ranked athletes benefit more).
This reconciles two competing predictions in the tournament literature:
\citet{GenakosPageliero2012} document that trailing athletes take more risks
when they face tougher opponents, while \citet{Brown2011} finds that top athletes
are discouraged by the presence of a superstar.
Our results suggest both effects are real but operate in different quality regimes:
\citeauthor{Brown2011}'s discouragement applies to top athletes in below-average
fields who face an unexpectedly strong rival, while \citeauthor{GenakosPageliero2012}'s
catch-up effect applies to athletes trailing in already-strong fields.

\paragraph{The team sports altruism confound.}
As argued in the Introduction, peer effects in team sports conflate social norms
with altruistic effort for teammates.
An athlete who accelerates because a teammate is struggling has been \textit{altruistically}
peer-affected---not because of conformism, tournament incentives, or rivalry, but
because of internalized preferences over the teammate's welfare.
Standard peer-effect regressions cannot distinguish this from incentive-based
peer effects.
Our individual-sport setting cleanly resolves this ambiguity: no mechanism
for altruistic effort exists.
The estimated peer effects are therefore pure measures of the social norm,
conformism, and rivalry channels.
This suggests that peer-effect estimates from team sports settings should be
interpreted with caution: the portion attributable to altruism is a
redistribution of effort within the team rather than an increase in total effort.

\paragraph{Implications for personnel economics.}
Our results suggest that standard tools of incentive design---prize spreads,
performance reviews, leaderboards---would have different effects on familiar versus
unfamiliar pairings of workers.
In workplaces where employees have long shared tenure (high familiarity), the
peer effect is larger, particularly for pairs where both are high performers.
Organizing workers into stable teams with known, capable peers would activate
the rivalry channel and generate the largest peer-effect multiplier.
Rotating workers across teams to prevent the formation of negative peer effects
(complacency against inferiors) would also be valuable based on our findings.

%% ====================================================================
\section{Conclusion}
\label{sec:conclusion}

We have documented three facts about peer effects in elite track athletics that
are novel to the economics literature.
First, the history of prior interaction between competitors---familiarity---is
a first-order determinant of the peer-effect magnitude, generating a factor-of-twelve
difference between strangers (one prior meeting) and rivals ($7+$ prior meetings).
Second, whether the peer effect is amplified or attenuated by familiarity depends
on the quality of familiar opponents: repeated competition against high-quality
rivals generates strong amplification, while repeated competition against inferiors
attenuates it.
Third, task complexity (race distance) moderates both the level and the familiarity
slope of peer effects, with middle-distance events exhibiting the strongest rivalry
dynamics consistent with an intermediate cost-precision regime.

Theoretically, we have developed a unified framework that nests all eight HOX
mechanisms and the familiarity channel in a single conformist-rivalry model.
The framework generates closed-form comparative statics that match all six
empirical predictions, and the proofs formalize the roles of the effort cost
function, information precision, rivalry intensity, and peer-statistic choice
in determining equilibrium effort.

More broadly, our results demonstrate the economic significance of rival
recognition as a channel of peer influence.
The standard assumption in the peer effects literature---that peers are
interchangeable given their observable characteristics---misses a quantitatively
important dimension: \textit{who} you have competed against before matters as
much as their quality level.
In individual-sport settings with observable performance and no altruism
confound, the rivalry channel can be cleanly identified.
The challenge for future work is to design identification strategies that
isolate the rivalry channel in settings with team production, strategic
interaction, and richer social relationships.

%% ====================================================================
\clearpage
\appendix

\section{Mathematical Proofs}
\label{app:proofs}

\paragraph{Proof of Proposition~\ref{prop:eight}.}

In a symmetric equilibrium where all athletes have the same quality $\theta$,
the FOC (\ref{eq:FOC_general}) becomes:
\begin{equation*}
W h_m(\theta) - c(\kappa) a^* - \frac{1}{\alpha}(a^* - \phi_m A_m^*) = 0
\end{equation*}
(rivalry term vanishes at symmetric equilibrium by normalization $V'(0)=0$).
Solving for $a^*$:
\begin{equation*}
a_m^* = \frac{W h_m(\theta) + \phi_m A_m^*/\alpha}{c(\kappa) + 1/\alpha}.
\end{equation*}
Differentiating with respect to $\theta$:
\begin{equation*}
\beta_m = \frac{\partial a_m^*}{\partial\theta}
= \frac{W \partial h_m/\partial\theta + \phi_m (\partial A_m^*/\partial\theta) /\alpha}
       {c(\kappa)+1/\alpha}.
\end{equation*}
Under the SOC (Assumption~\ref{ass:reg}(ii)), the denominator is positive.
$\square$

\paragraph{Proof of Corollary~\ref{cor:ranking}.}

For Shining Light: $A^* = \max_j s_j^- = \theta + \delta$, where $\delta$ is a
positive shift due to within-field heterogeneity.
$\partial A_{\text{Shining}}^*/\partial\theta = 1$ exactly.
The marginal winning probability $h_{\text{Shining}} = \partial P/\partial a$ is also higher
under Shining Light than under LIM, because winning requires beating the fastest
opponent---whose quality level is more binding than the mean \citep{LazearRosen1981}.

For Bad Apple: $A^* = \min_j s_j^- = \theta - \delta'$.
$\partial A_{\text{Bad Apple}}^*/\partial\theta < 1$ (the minimum shifts less than the mean
when heterogeneity is present, and in the limit for large $n$,
$\partial \min_j / \partial\theta \to 0$ since the minimum is determined
by tail behavior).
Moreover, $h_{\text{Bad Apple}} \approx 0$: in a tournament, winning probability
depends on beating the best opponents, not the worst.
Hence $\beta_{\text{Bad Apple}} \approx 0$.
$\square$

\paragraph{Proof of Proposition~\ref{prop:fam}.}

Substitute $\phi_i(\bar{f}) = \phi_0 + \phi_1 \bar{f}$ and the familiarity-weighted
peer statistic into the FOC and differentiate using the implicit function theorem.
The key steps are:
\begin{align*}
\frac{\partial a^*}{\partial\theta} &=
\frac{W\partial h/\partial\theta + [\phi_0+\phi_1\bar{f}]\partial A^*/\partial\theta/\alpha
+ \gamma\bar{f} V''(0)\partial\mu^f/\partial\theta}
{c(\kappa)+1/\alpha - \gamma\bar{f} V''(0)}.
\end{align*}
The numerator is increasing in $\bar{f}$ when $\phi_1 > 0$;
the denominator is decreasing in $\bar{f}$ when $\gamma V''(0) > 0$.
Both effects amplify $\beta_m(\bar{f})$.
$\square$

\paragraph{Proof of Proposition~\ref{prop:complexity}.}

From (\ref{eq:beta_fam}):
$\frac{\partial\beta_m}{\partial\kappa}
= \frac{[\phi_0+\phi_1\bar{f}]\phi'(\kappa)\partial A^*/\partial\theta/\alpha \cdot D
- N \cdot c'(\kappa)}{D^2}$,
where $N$ is the numerator and $D = c(\kappa)+1/\alpha-\gamma\bar{f} V''(0)$.
For $\kappa < \kappa^*$: $\phi'(\kappa) > 0$ (rising) and $c'(\kappa) < 0$ (falling),
so both terms in the numerator are positive $\Rightarrow$ $\partial\beta_m/\partial\kappa > 0$.
For $\kappa > \kappa^*$: $\phi'(\kappa) < 0$ and $c'(\kappa) > 0$, so
$\partial\beta_m/\partial\kappa < 0$.
$\square$

\section{Additional Tables}
\label{app:tables}

\begin{table}[H]
\centering
\begin{threeparttable}
\caption{Quality-Conditional Peer Effects by Familiarity and HOX Mechanism}
\label{tab:qual_full}
\small
\begin{tabular}{llcccccccc}
\toprule
& & \multicolumn{4}{c}{High-quality field ($z_{\bar{s}^-} > 0$)}
  & \multicolumn{4}{c}{Low-quality field ($z_{\bar{s}^-} \leq 0$)} \\
\cmidrule(lr){3-6}\cmidrule(lr){7-10}
HOX & Fam. bin & $\hat\beta$ & SE & Sig. & $N$ & $\hat\beta$ & SE & Sig. & $N$ \\
\midrule
\multirow{4}{*}{\parbox{2.5cm}{Shining Light}}
& Never    & & & & & & & & \\
& $1\times$& & & & & & & & \\
& $2$--$3\times$ & & & & & & & & \\
& $7+\times$ & & & & & & & & \\
\midrule
\multirow{4}{*}{\parbox{2.5cm}{Bad Apple}}
& Never    & & & & & & & & \\
& $1\times$ & & & & & & & & \\
& $2$--$3\times$ & & & & & & & & \\
& $7+\times$ & & & & & & & & \\
\bottomrule
\end{tabular}
\begin{tablenotes}\small
\item \textit{Notes:} Coefficients from table3\_hox\_quality\_familiarity.csv
(generated by peer\_analysis\_v2.py). Values populated from extended analysis.
\end{tablenotes}
\end{threeparttable}
\end{table}

%% ====================================================================
\clearpage
\singlespacing
\bibliographystyle{apalike}
\begin{thebibliography}{99}

\bibitem[Ager et al.(2022)]{AgerBursztynLeuchtVoth2022}
Ager, P., Bursztyn, L., Leucht, L., and Voth, H.-J. (2022).
\newblock Killer incentives: rivalry, performance and risk-taking among German
fighter pilots, 1939--45.
\newblock {\em Review of Economic Studies}, 89(5):2257--2292.

\bibitem[Angrist(2014)]{Angrist2014}
Angrist, J.~D. (2014).
\newblock The perils of peer effects.
\newblock {\em Labour Economics}, 30:98--108.

\bibitem[Arag\'on and Rezazadeh(2025)]{AragonRezazadeh2025}
Arag\'on, F.~M. and Rezazadeh, M. (2025).
\newblock Not too hot to run: adaptation to extreme heat in competitive runners.
\newblock Working Paper, Simon Fraser University.

\bibitem[Arcidiacono et al.(2017)]{Arcidiacono2017}
Arcidiacono, P., Kinsler, J., and Price, J. (2017).
\newblock Productivity spillovers in team production: evidence from
professional basketball.
\newblock {\em Journal of Labor Economics}, 35(1):191--225.

\bibitem[Ashby(2025)]{Ashby2025}
Ashby, N.~J. (2025).
\newblock Heterogeneous peer effects in competitive environments: the case of
high-school harriers.
\newblock {\em Applied Economics Letters}, 32(4):557--562.

\bibitem[Ashby and Ramos(2023)]{AshbyRamos2023}
Ashby, N.~J. and Ramos, A. (2023).
\newblock Cross-country runners and strategic shirking.
\newblock {\em Journal of Behavioral and Experimental Economics}, 104:101976.

\bibitem[B\'enabou and Tirole(2006)]{BenabouTirole2006}
B\'enabou, R. and Tirole, J. (2006).
\newblock Incentives and prosocial behavior.
\newblock {\em American Economic Review}, 96(5):1652--1678.

\bibitem[B\'enabou and Tirole(2011)]{BenabouTirole2011}
B\'enabou, R. and Tirole, J. (2011).
\newblock Laws and norms.
\newblock NBER Working Paper No. 17579.

\bibitem[Bernheim(1994)]{Bernheim1994}
Bernheim, B.~D. (1994).
\newblock A theory of conformity.
\newblock {\em Journal of Political Economy}, 102(5):841--877.

\bibitem[Boucher et al.(2024)]{Boucher2024}
Boucher, V., Rendall, M., Ushchev, P., and Zenou, Y. (2024).
\newblock Toward a general theory of peer effects.
\newblock {\em Econometrica}, 92(2):543--565.

\bibitem[Brown(2011)]{Brown2011}
Brown, J. (2011).
\newblock Quitters never win: the (adverse) incentive effects of competing with superstars.
\newblock {\em Journal of Political Economy}, 119(5):982--1013.

\bibitem[Cohen-Zada et al.(2024)]{CohenZada2024}
Cohen-Zada, D., Dayag, T., and Gershoni, N. (2024).
\newblock Peer effects in competitive team sports: evidence from professional soccer.
\newblock {\em Management Science}.

\bibitem[Cornelissen et al.(2017)]{CornelissenDustmannSchonberg2017}
Cornelissen, T., Dustmann, C., and Sch\"{o}nberg, U. (2017).
\newblock Peer effects in the workplace.
\newblock {\em American Economic Review}, 107(2):425--456.

\bibitem[Emerson and Hill(2018)]{EmersonHill2018}
Emerson, J. and Hill, B. (2018).
\newblock Peer effects in marathon running: evidence from individual performance data.
\newblock {\em Labour Economics}, 55:170--180.

\bibitem[Falk and Ichino(2006)]{FalkIchino2006}
Falk, A. and Ichino, A. (2006).
\newblock Clean evidence on peer effects.
\newblock {\em Journal of Labor Economics}, 24(1):39--57.

\bibitem[Genakos and Pagliero(2012)]{GenakosPageliero2012}
Genakos, C. and Pagliero, M. (2012).
\newblock Interim rank, risk taking, and performance in dynamic tournaments.
\newblock {\em Journal of Political Economy}, 120(4):782--813.

\bibitem[Gould and Winter(2009)]{GouldWinter2009}
Gould, E.~D. and Winter, E. (2009).
\newblock Interactions between workers and the technology of production:
evidence from professional baseball.
\newblock {\em Review of Economics and Statistics}, 91(1):188--200.

\bibitem[Green and Stokey(1983)]{GreenStokey1983}
Green, J.~R. and Stokey, N.~L. (1983).
\newblock A comparison of tournaments and contracts.
\newblock {\em Journal of Political Economy}, 91(3):349--364.

\bibitem[Guryan et al.(2009)]{GuryanKroftNotowidigdo2009}
Guryan, J., Kroft, K., and Notowidigdo, M.~J. (2009).
\newblock Peer effects in the workplace: evidence from random groupings in professional golf tournaments.
\newblock {\em American Economic Journal: Applied Economics}, 1(4):34--68.

\bibitem[Hill(2014)]{Hill2014}
Hill, B. (2014).
\newblock The heat is on: tournament structure, peer effects, and performance.
\newblock {\em Journal of Sports Economics}, 15(4):315--337.

\bibitem[Hoxby and Weingarth(2007)]{HoxbyWeingarth2007}
Hoxby, C.~M. and Weingarth, G. (2007).
\newblock Taking race out of the equation: school reassignment and the
structure of peer effects.
\newblock Working Paper, Harvard University.

\bibitem[Jiang(2020)]{Jiang2020}
Jiang, Y. (2020).
\newblock Peer effects in elite swimming.
\newblock {\em Journal of Economic Behavior \& Organization}, 171:35--54.

\bibitem[Kilduff(2014)]{Kilduff2014}
Kilduff, G.~J. (2014).
\newblock Driven to win: rivalry, motivation, and performance.
\newblock {\em Social Psychological and Personality Science}, 5(8):944--952.

\bibitem[Kilduff et al.(2010)]{KilduffElfenbeinStaw2010}
Kilduff, G.~J., Elfenbein, H.~A., and Staw, B.~M. (2010).
\newblock The psychology of rivalry: a relationally dependent analysis of competition.
\newblock {\em Academy of Management Journal}, 53(5):943--969.

\bibitem[Lazear and Rosen(1981)]{LazearRosen1981}
Lazear, E.~P. and Rosen, S. (1981).
\newblock Rank-order tournaments as optimum labor contracts.
\newblock {\em Journal of Political Economy}, 89(5):841--864.

\bibitem[Manski(1993)]{Manski1993}
Manski, C.~F. (1993).
\newblock Identification of endogenous social effects: the reflection problem.
\newblock {\em Review of Economic Studies}, 60(3):531--542.

\bibitem[Mas and Moretti(2009)]{MasMoretti2009}
Mas, A. and Moretti, E. (2009).
\newblock Peers at work.
\newblock {\em American Economic Review}, 99(1):112--145.

\bibitem[Molodchik et al.(2021)]{MoloChik2021}
Molodchik, M., Paklina, S., and Parshakov, P. (2021).
\newblock Peer effects in a team production setting: evidence from soccer.
\newblock {\em German Economic Review}, 22(4):449--468.

\bibitem[Nalebuff and Stiglitz(1983)]{NalebuffStiglitz1983}
Nalebuff, B.~J. and Stiglitz, J.~E. (1983).
\newblock Prizes and incentives: towards a general theory of compensation and competition.
\newblock {\em Bell Journal of Economics}, 14(1):21--43.

\bibitem[Pike et al.(2018)]{PikeKilduffGalinsky2018}
Pike, B.~E., Kilduff, G.~J., and Galinsky, A.~D. (2018).
\newblock The long shadow of rivalry: rivalry motivates performance today and tomorrow.
\newblock {\em Psychological Science}, 29(5):804--813.

\bibitem[Sacerdote(2001)]{Sacerdote2001}
Sacerdote, B. (2001).
\newblock Peer effects with random assignment: results for Dartmouth roommates.
\newblock {\em Quarterly Journal of Economics}, 116(2):681--704.

\bibitem[Yamane and Hayashi(2015)]{YamaneHayashi2015}
Yamane, S. and Hayashi, R. (2015).
\newblock Peer effects among swimmers.
\newblock {\em Scandinavian Journal of Economics}, 117(4):1230--1255.

\end{thebibliography}

\end{document}
